\section{Introducción}

\subsection{Contexto}

En un colapso estructural, el tiempo decide. En Kobe (1995) salió con vida cerca del
80~\% de las personas liberadas el primer día, cerca del 30~\% en los días 2 y 3, y casi
nadie a partir del cuarto~\cite{kawata}. En Italia (1980), el 93~\% de los atrapados que
sobrevivieron fue sacado en las primeras 24~horas~\cite{noji}. La regla de las «72 horas»
resume ese decaimiento, aunque no existe un estudio que la derive y hay rescates
documentados hasta 14 días después~\cite{macintyre}.

A la mayoría la sacan sus vecinos: en Kobe, cerca del 80~\% de los rescatados de edificios
colapsados fue liberado por familiares y vecinos~\cite{japon2014}. Los equipos de rescate
y su tecnología trabajan sobre la fracción difícil, la de quienes quedaron en huecos
profundos a los que no se llega con las manos.

El sismo del 19 de septiembre de 2017 (magnitud 7.1, epicentro en el límite entre Puebla
y Morelos~\cite{ssn2017}) dejó 369 muertos, 45 de ellos en el estado de Puebla, y 44
edificios colapsados en la Ciudad de México (otros conteos dan 38 o 46)~\cite{galvis2020}.
El 91~\% de esos edificios era anterior a 1985. Puebla volverá a necesitar esta capacidad.

\subsection{El problema}

Las herramientas actuales no entran donde está la víctima (tabla~\ref{tab:tecnologias}).
El perro trabaja la superficie; la cámara en poste alcanza unos 5~m en línea casi recta; el
robot más pequeño usado en campo necesita un hueco de $6\times17$~cm. En el World Trade
Center los huecos explorados eran, por lo general, de menos de 1~m~\cite{murphy2004}, y
hasta 2021 ningún robot terrestre había encontrado a una persona viva en un
desastre~\cite{murphy2021}. El robot serpiente que Carnegie Mellon llevó a la Ciudad de
México en 2017 (5~cm de diámetro) entró a un edificio, hizo dos pasadas y no encontró a
nadie~\cite{whitman2018}.

Una cucaracha de Madagascar mide 6~cm de largo y 1.3~cm de alto, trepa obstáculos de su
propia altura y sigue paredes sin que nadie se lo pida. Varios grupos la han convertido
en plataforma de sensores, y en 2025 diez de ellas se usaron por primera vez en un
desastre real, en Myanmar, sin hallar sobrevivientes y con fallas técnicas
reconocidas~\cite{htx2025}. Lo que falta no es la idea, sino un sistema barato, sin cirugía,
que no requiera pilotear a cada insecto, cuyo circuito quepa en lo que el animal puede cargar
y que le diga al rescatista dónde está la persona aunque el insecto que la encontró haya
quedado fuera del alcance de la radio.

\subsection{Pregunta, hipótesis y objetivos}

\begin{infobox}
\textbf{Pregunta de investigación.} ¿Cuánto mejora la búsqueda de una persona atrapada si
la mochila, además de detectar, guía al insecto con estímulos mínimos; qué circuito digital
cabe en la carga que el insecto soporta, y cómo se coordinan varias mochilas para ubicar a
la persona sin que una computadora haga las cuentas?
\end{infobox}

\textbf{Hipótesis.} Una política que corre a bordo y estimula poco (solo cuando el insecto
se detiene, de vez en cuando para despegarlo de las paredes y para acercarlo a una mancha
de calor), con un presupuesto que limita los estímulos por hora, aumenta la probabilidad de
detección en una sesión de 45 a 60~minutos respecto a la exploración pasiva, sin provocar
habituación, y cabe en una FPGA de miligramos. Dos mochilas que detectan a la misma persona
bastan para que un receptor, también en FPGA, estime su posición cruzando los dos rumbos.

\textbf{Objetivo general.} Diseñar, describir en Verilog, verificar y probar en la DE10-Lite
el sistema digital de una mochila de búsqueda para \Gp{} y la red con la que varias mochilas
localizan a una persona y avisan a un operador; evaluar por simulación cuánto aporta, y
llevar el núcleo a un layout en un proceso abierto.

\textbf{Objetivos específicos.}
\begin{enumerate}
  \item Establecer, con datos publicados, qué carga, velocidad y terreno soporta el insecto.
  \item Comparar los tipos de estímulo y modelar la respuesta y la habituación.
  \item Dimensionar la mochila: masa, energía, costo y capacidad de cómputo.
  \item Definir un método de detección y de localización que funcione a distancia.
  \item Diseñar la política de estímulos, describirla en Verilog, verificarla contra un modelo
        de referencia y sintetizarla.
  \item Diseñar una red de tres nodos en FPGA: un protocolo de enlace con detección de errores y
        retransmisión, el reconocimiento de vecinos por tonos y el cálculo de distancias y de la
        posición de la persona en punto fijo, de modo que la computadora solo muestre resultados.
  \item Verificar la red con bancos de prueba autoverificables, incluidos los casos límite y de
        error, y medir recursos, frecuencia y latencia.
  \item Simular la misión completa y medir tiempo a la detección, cobertura y estímulos aplicados.
  \item Llevar el núcleo a un layout en SKY130 con un flujo abierto y compararlo con la FPGA.
  \item Analizar la practicidad, el costo y las implicaciones éticas y legales.
\end{enumerate}

\subsection{Contribuciones}

\begin{itemize}
  \item \textbf{Un núcleo que cierra el lazo a bordo.} Detección, decisión y estímulo en 477 elementos
        lógicos, idéntico a su modelo de referencia en 11\,424 cuadros y probado en dos tarjetas reales.
  \item \textbf{Una red de tres FPGA en la que la computadora no calcula nada.} Tramas con CRC-16 y
        número de secuencia, retransmisión por una mochila puente, reconocimiento de vecinos con
        filtros de Goertzel, distancias exactas y triangulación con CORDIC. Nueve escenarios de
        prueba, con ruido en la línea y con el corte del alcance.
  \item \textbf{Aritmética en punto fijo con el error medido.} El error del CORDIC según el número de
        iteraciones y la comparación que llevó a calcular las distancias con una raíz cuadrada
        entera en lugar del CORDIC.
  \item \textbf{El núcleo como circuito integrado.} Un layout en SKY130 de 0.050~mm$^2$, limpio en DRC
        y LVS, que sustituye a la estimación por conteo de celdas.
  \item \textbf{Una evaluación de la misión con sus límites.} \numMisiones{} misiones simuladas por caso,
        análisis de sensibilidad y el caso en que guiar empeora el resultado.
  \item \textbf{Un repositorio que reproduce cada resultado:} código, bancos de prueba, proyectos,
        layout y los programas de las figuras (apéndice~\ref{sec:repositorio}).
\end{itemize}

\subsection{Alcance}

\begin{itemize}
  \item \textbf{Sí incluye:} el modelo completo, la simulación de misión, el núcleo y la red
        verificados y sintetizados, las pruebas en una, dos y tres DE10-Lite, y el layout del núcleo.
  \item \textbf{Simulado dentro de la FPGA:} la matriz térmica (los interruptores fijan la temperatura
        que lee), el canal acústico de los tonos y el radio, que en la demostración es un cable. La
        posición de cada mochila se da por conocida.
  \item \textbf{No incluye:} cirugía ni electrodos implantados, pruebas con insectos, pruebas en
        escombro real ni la fabricación del chip.
  \item \textbf{Evidencia:} todo número proviene de una fuente citada, de la síntesis en Quartus,
        de la simulación en Icarus Verilog, del flujo de layout o de la simulación de misión; los
        supuestos se marcan como tales.
\end{itemize}

El documento sigue este orden: estado del arte (sección~\ref{sec:arte}), metodología
(sección~\ref{sec:metodologia}), cuatro casos de estudio (sección~\ref{sec:casos}), análisis
comparativo (sección~\ref{sec:comparativo}), ética e impacto (sección~\ref{sec:etica}) y
conclusiones con trabajo futuro (sección~\ref{sec:conclusiones}).
